\begin{summarybox}
	\textbf{\textcolor{CSummary}{一句话核心}}

	加减法看各项是否一起变号或一起不变；
	乘除法看变号次数。判断前先检查运算结果的定义域。

	\medskip
	\textbf{\textcolor{CSummary}{规则速查}}
	\[
		\begin{aligned}
			\text{奇}\pm\text{奇}    & =\text{奇},
			                       & \qquad \text{偶}\pm\text{偶} & =\text{偶},   \\
			\text{奇}\times\text{奇} & =\text{偶},
			                       & \text{奇}\times\text{偶}     & =\text{奇},   \\
			\text{偶}\times\text{偶} & =\text{偶}.                 &            &
		\end{aligned}
	\]
	除法的奇偶规则与乘法相同，但商只在分母不为零的点上定义。

	\medskip
	\textbf{\textcolor{CSummary}{两个提醒}}

	\begin{itemize}[leftmargin=*, itemsep=0.4em, topsep=0.3em]
		\item 奇函数与偶函数相加减，通常既非奇也非偶；
		      零函数等特殊情况要单独判断。
		\item 除法删去分母的零点后，仍要保留原式的定义域限制；
		      不能因约分而把删去的点加回来。
	\end{itemize}
\end{summarybox}